% ============================================================
\section{Adeguatezza}
\label{sec:adeguatezza}
% ============================================================

Due metriche di adeguatezza, entrambe misurate con JaCoCo: branch decision coverage e statement
coverage (nell'output di JaCoCo quest'ultima è riportata per riga, colonna \texttt{LINE}).

\subsection{BrokerImpl}

Valori calcolati sull'insieme di test Maven (BB, Randoop, LLM;
\hyperref[cmd:jacoco-brokerimpl]{Appendice}): branch decision coverage 17,69\% (253/1430),
statement coverage 27,07\% (680/2512). Il divario di quasi 10 punti conferma la gerarchia: molte
istruzioni vengono eseguite senza che siano stati esercitati entrambi i rami di ogni decisione al
loro interno.

\texttt{BrokerImpl} è una classe di 5510 righe e circa 250 metodi pubblici o protetti: nessuna
delle suite ha mai avuto l'obiettivo di coprirla per intero. La suite BB copre per progettazione
solo i 7 metodi selezionati nel Capitolo~\ref{sec:cp}, non un campione rappresentativo dell'intera
classe.

Rilevanza rispetto alle funzionalità attese: i test BB derivano l'oracolo da Javadoc, quindi ogni
ramo o riga che coprono corrisponde per costruzione a comportamento documentato. Randoop genera
senza conoscere la specifica: aumenta la coverage ma non garantisce di esercitare comportamento
rilevante secondo la documentazione della classe. I test LLM occupano una posizione intermedia:
hanno accesso al codice sorgente nel prompt, ma producono asserzioni motivate da un ragionamento
sul comportamento atteso, non dedotte meccanicamente dai rami del codice. \texttt{EvoSuite}, mai
integrato nella build Maven, non contribuisce a questa misurazione.

\subsubsection{Coverage (JaCoCo)}

\textit{Selezione.}
\texttt{lock(Object, OpCallbacks)} e \texttt{detachAll(OpCallbacks, boolean)} risultano a 0\%,
accanto ai rispettivi overload già testati nel Capitolo~\ref{sec:cp},
\texttt{lock(Object, int, int, OpCallbacks)} all'87\% e \texttt{detachAll(Collection,
OpCallbacks)} al 100\%. La firma dei due metodi, già nota dalla Javadoc, basta a riconoscerli come
metodi pubblici mai esercitati: nessuna lettura del corpo necessaria. Tra gli altri metodi
segnalati da JaCoCo la maggioranza sono helper interni poco significativi; selezionati invece
\texttt{delete(Object, OpCallbacks)} e \texttt{flush()} per rilevanza funzionale.

\textit{Progettazione.} \texttt{lock(Object pc, OpCallbacks call)}: riusa le categorie di
\texttt{pc} e il comportamento di \texttt{call.processArgument} già stabiliti per l'overload a 4
argomenti. A differenza di quello, \texttt{level} e \texttt{timeout} non sono parametri: la
Javadoc dichiara il livello corrente della \texttt{FetchConfiguration}, il timeout non è
documentato. Ipotesi: \texttt{ACT\_NONE}/\texttt{ACT\_CASCADE}/\texttt{ACT\_RUN} equivalenti tra
loro; per \texttt{pc} non gestito, nessun effetto e nessuna eccezione, per analogia col risultato
già confermato per l'overload a 4 argomenti su \texttt{pc} non gestito nel
Capitolo~\ref{sec:implementazione}. Classi in Tabella~\ref{tab:lock2-classi}. 20 tuple.

\texttt{detachAll(OpCallbacks call, boolean flush)}: non ha un parametro con gli oggetti da
staccare, opera su tutte le istanze già tracciate dal broker. Un'istanza entra in quel registro
solo chiamando prima un altro metodo pubblico, tipicamente \texttt{persist()}: non essendoci un
parametro da cui prenderle, l'unica via è costruire lo stato del broker prima di chiamare
\texttt{detachAll}. Categorie: \texttt{call} (null/non null) $\times$ \texttt{flush} $\times$
transazione $\times$ stato del broker, \{registro vuoto\} oppure \{registro con un'istanza
gestita\}, dove "gestita" vuol dire semplicemente presente in quel registro, indipendentemente dal
fatto che sia nuova o già presente nello store. Con registro vuoto, transazione incrociata con
\texttt{call} e \texttt{flush}: 8 combinazioni. Con registro con un'istanza gestita, transazione
fissa ad attiva, richiesta da \texttt{persist()} per costruire quello stato: solo \texttt{call}
$\times$ \texttt{flush}, altre 4 combinazioni, 12 in totale. Ipotesi per il registro vuoto:
nessuna eccezione, callback non invocato per assenza di istanze da processare. Ipotesi per il
registro con un'istanza gestita, per analogia con \texttt{detachAll(Collection, OpCallbacks)} del
Capitolo~\ref{sec:cp}: nessuna eccezione, callback invocato sull'istanza gestita. Classi in
Tabella~\ref{tab:detachall2-classi}. 12 tuple.

\texttt{delete(Object pc, OpCallbacks call)}: Javadoc minima, "Delete the given object.". Stesse
categorie di \texttt{pc} e stesso comportamento di \texttt{call.processArgument} di \texttt{lock}.
Né \texttt{pc} null né \texttt{pc} non gestito sono specificati dalla Javadoc: ipotesi per
entrambi, nessun effetto, per analogia con gli altri metodi della classe che si comportano così su
un oggetto non ancora persistito. Classi in Tabella~\ref{tab:delete-classi}. 20 tuple.

\texttt{flush()}: nessun parametro formale, categorie dallo stato del broker. \texttt{transazione}:
\{non attiva\}, \{attiva\}. \texttt{istanze da flushare}: \{nessuna\}, \{presenti\}.
\texttt{esito dello store}* (richiede istanze presenti, senza scrittura non c'è un esito da
misurare): \{successo\}, \{fallimento\}, quest'ultimo documentato dalla Javadoc, "may set the
rollback only flag... if it encounters an error". Prodotto cartesiano: $2 \times 2 \times 2 = 8$
combinazioni. Con \{nessuna\} istanza l'esito è privo di senso, non c'è una scrittura da cui
misurarlo: le combinazioni \{non attiva, nessuna istanza, successo\} e \{non attiva, nessuna
istanza, fallimento\} sono indistinguibili e si uniscono in una sola, \{non attiva, nessuna
istanza\}; allo stesso modo \{attiva, nessuna istanza, successo\} e \{attiva, nessuna istanza,
fallimento\} si uniscono in \{attiva, nessuna istanza\}. Dalle 8 combinazioni di partenza ne
restano 6. Per ipotesi dalla Javadoc: senza transazione attiva non ci sono comunque istanze
da flushare, quindi le combinazioni \{non attiva, istanze presenti, successo\} e \{non attiva,
istanze presenti, fallimento\} non sono modellate, escluse per ipotesi e non per impossibilità
accertata. Restano 4 categorie progettate: \{non attiva\},
\{attiva, nessuna istanza\}, \{attiva, istanze presenti, successo\}, \{attiva, istanze presenti,
fallimento\}. Ipotesi sull'oracolo: nessuna eccezione nei primi tre casi; \texttt{rollbackOnly} a
true nel quarto. Classi in Tabella~\ref{tab:flush-classi}. 4 categorie.

\textit{Implementazione e correzione.} \texttt{lock}: l'ipotesi iniziale, i tre valori di azione
equivalenti tra loro, non regge per \texttt{pc} gestito: solo \texttt{ACT\_RUN} fa procedere il
lock, \texttt{ACT\_NONE}/\texttt{ACT\_CASCADE} non hanno effetto. Per \texttt{pc} null e non
gestito l'equivalenza regge, perché nessuno dei tre ha effetto in nessun caso. Ipotesi su
\texttt{pc} non gestito confermata. 20/20 test passano.

\texttt{detachAll}: con broker senza istanze gestite, confermato con \texttt{verify(callback,
never()).processArgument(...)} che il callback non viene mai invocato, in tutte le 8 tuple. Con
un'istanza gestita, ottenuta persistendo davvero un'istanza con la stessa tecnica di registrazione
usata in \texttt{PersistTest} del Capitolo~\ref{sec:implementazione}: l'ipotesi di progettazione,
nessuna eccezione e callback invocato, è sbagliata. Con \texttt{flush=false} l'esecuzione rivela
\texttt{InvalidStateException}, un'istanza appena persistita e mai flushata non può essere
staccata; il callback non viene invocato affatto, verificato con \texttt{verify(callback,
never()).processArgument(...)}. Oracolo corretto di conseguenza a eccezione attesa: le 2 tuple con
\texttt{flush=false}, \texttt{call} null e non null, sono implementate. Con \texttt{flush=true}
il test fallisce con un'eccezione il cui stack trace, output dell'esecuzione, non lettura del
sorgente, indica \texttt{StateManagerImpl.assertObjectIdAssigned}: richiede un'assegnazione reale
dell'id dallo store, che il mock minimale usato qui non fornisce. Le 2 tuple con \texttt{flush=true}
restano non implementate. In totale 10/12 tuple implementate, tutte passano: le 8 con broker senza
istanze gestite più le 2 appena corrette con \texttt{flush=false}.

\texttt{delete}: nessuna eccezione confermata per \texttt{pc} gestito, indipendentemente
dall'azione. Per \texttt{pc} non gestito, ipotesi per analogia confermata per \texttt{call}
assente, \texttt{ACT\_NONE}, \texttt{ACT\_RUN}. \texttt{ACT\_CASCADE} inizialmente fallisce con
un'eccezione: lo stack trace dell'esecuzione indica \texttt{cascadeTransient}, che tratta l'istanza
non gestita come se fosse nuova e prova a persisterla, creando un vero \texttt{StateManager} che
richiede il tipo di identità della classe, dichiarato nella \texttt{ClassMetaData}. Il mock di
\texttt{pc} non gestito usato per le altre azioni non ha questa informazione: aggiunta registrando
la classe con \texttt{PCRegistry.register} e collegandole una \texttt{ClassMetaData} mock con
\texttt{getIdentityType()} impostato, stessa configurazione già scritta per il caso \texttt{pc} non
gestito di \texttt{persist()} in \texttt{PersistTest} del Capitolo~\ref{sec:implementazione}, qui
copiata sull'istanza di \texttt{delete}. Con questa aggiunta l'esecuzione conferma l'ipotesi anche
per \texttt{ACT\_CASCADE}, nessuna eccezione. 20/20 tuple implementate, tutte passano.

\texttt{flush}: ipotesi confermata per i primi due casi, nessuna correzione. Caso di successo:
stesso muro di \texttt{StateManagerImpl.assertObjectIdAssigned} già incontrato per
\texttt{detachAll(flush=true)}; implementazione non proseguita.

Caso di fallimento. \texttt{BrokerImpl.flush()} non scrive lui stesso sul datastore: delega a uno
\texttt{StoreManager}, mockato al posto di un database reale. La sua firma dichiarata,
\texttt{Collection<Exception> flush(Collection<OpenJPAStateManager> sms)}, prova a scrivere tutte
le istanze ricevute e restituisce le eccezioni delle sole istanze fallite, invece di lanciarne una
sola e fermarsi: lista vuota, tutto scritto; lista non vuota, quelle istanze non si sono scritte.
Simulato un fallimento facendo restituire a \texttt{storeManager.flush()} una
lista con un'eccezione. Ipotesi, da Javadoc: \texttt{BrokerImpl}, vedendo quella lista non vuota,
imposta \texttt{rollbackOnly} a true. Osservato invece: \texttt{rollbackOnly} resta false, nessuna
eccezione lanciata, test senza errori. A differenza del caso di successo, qui non c'è uno stack
trace che indichi cosa manca al mock: non è possibile stabilire se \texttt{BrokerImpl} abbia davvero ispezionato
quella lista in questo scenario o l'abbia ignorata, e senza un errore da diagnosticare non c'è modo
di saperlo senza leggere il codice. Implementazione non proseguita. In totale 2/4 categorie
implementate, i 2 test corrispondenti passano.

Risultato: branch decision coverage 20,91\% (299/1430), statement coverage 30,25\% (760/2512), da
17,69\% (253/1430) e 27,07\% (680/2512) di partenza: +3,22pp branch, +3,18pp statement. Questo
numero finale è misurato sulla suite dopo entrambe le iterazioni: include quindi anche il piccolo
contributo di \texttt{delete2}/\texttt{newinstance2} (sottosezione successiva), selezionati
dall'analisi dei mutanti sopravvissuti e non da questa misura di copertura.

Un'analisi white-box dei rami scoperti alzerebbe questi numeri più rapidamente, ma sposterebbe
l'obiettivo dalla verifica delle funzionalità documentate al solo inseguimento della copertura:
il primo resta l'obiettivo principale di questo progetto. Analisi condotta \textit{best effort}
sulla documentazione disponibile, non sul codice.

\subsubsection{Mutation testing (PITest)}

\textit{Selezione.} Mutazioni sull'intera classe, dalla suite dopo l'iterazione di coverage: 1475
generate, 259 uccise (17,56\%). PITest aggrega per nome di metodo, non per overload: le mutazioni
di \texttt{lock}, \texttt{detachAll} e \texttt{find} includono anche overload mai selezionati nel
Capitolo~\ref{sec:cp}, ad esempio \texttt{find(Object, FetchConfiguration, BitSet, ...)} contro
l'unico testato, \texttt{find(Object, boolean, FindCallbacks)}. Su questi tre nomi, la
maggioranza dei mutanti non uccisi risulta \texttt{NO\_COVERAGE} sull'overload mai selezionato,
non \texttt{SURVIVED} sull'overload progettato. Per \texttt{attachAll}, \texttt{isDetached},
\texttt{find}, \texttt{lock}, \texttt{detachAll}, \texttt{persist} la progettazione copre già
ogni dimensione dichiarata dalla Javadoc, sull'overload testato: margine minimo, \texttt{persist}
a punteggio pieno. \texttt{delete}, 12 sopravvissuti su 23, e \texttt{newInstance}, 3 su 7,
partivano da una progettazione meno matura: selezionati per il miglioramento. L'overload non
selezionato di \texttt{find} richiederebbe una progettazione Category Partition nuova per una
firma mai analizzata: non è stato fatto.

\textit{Progettazione.} \texttt{delete}: nuova categoria per \texttt{pc}, gestito e detached,
riusata dalla categoria "detached" già stabilita dalla Javadoc di \texttt{isDetached()} del
Capitolo~\ref{sec:cp}. Ipotesi: eccezione attesa, tipo non specificato dal design. Classi in
Tabella~\ref{tab:delete2-classi}. 2 tuple.

\texttt{newInstance(Class cls)}: la Javadoc dichiara \texttt{@throws IllegalArgumentException if
cls is not a managed type or interface}, oracolo più preciso di quello già usato per il caso "tipo
non managed concreto". Classi in Tabella~\ref{tab:newinstance2-classi}. 1 tupla.

\textit{Implementazione e correzione.} \texttt{delete}: ipotesi confermata, \texttt{UserException},
"You cannot perform operation delete on detached object...". 2/2 test passano. Mutation score
finale: 10/23 mutanti uccisi, da 6/23 di partenza; il miglioramento viene sia dai 2 test qui
progettati sia dalle altre tuple \texttt{pc} non gestito già aggiunte per \texttt{delete}.

\texttt{newInstance}: oracolo confermato direttamente da Javadoc, nessuna correzione; 4/7
invariato, nessun mutante nuovo raggiunto.

Risultato: punteggio aggregato sull'intera classe 272/1475 (18,44\%), da 259/1475 (17,56\%) di
partenza.

Stesso compromesso della coverage: leggere il diff dei mutanti sopravvissuti alzerebbe il
punteggio più rapidamente, a scapito della verifica delle funzionalità documentate. Anche questa
analisi resta \textit{best effort} sulla documentazione, non sul codice.

Tuple complete per tutti i metodi in \texttt{report/data/}.

\subsubsection{Reliability}

Assumendo un profilo operazionale con probabilità di utilizzo
uniforme sull'insieme finale di test ottenuto dopo le iterazioni dei paragrafi precedenti:
reliability $= 1 - PFD =$ test passati / test eseguiti.

La suite Maven per \texttt{BrokerImpl} conta 623 test, esecuzioni di metodi \texttt{@Test}
incluse le ripetizioni parametrizzate di \texttt{@Parameters}. 15 sono disabilitati con
\texttt{@Ignore} e non fanno parte dell'insieme di esecuzioni misurate: 3 nei test LLM
(\texttt{BrokerImplLlmZ1Test}, \texttt{BrokerImplLlmZ2Test}) hanno un oracolo scorretto, non un
comportamento scorretto della classe; 12 in \texttt{RandoopErrorTest} sono sequenze
generate su un'istanza mai inizializzata.

Sui restanti 608 test eseguiti: 608 passano, 0 falliscono. Reliability $=$ 608/608 $=$ 1
(PFD $=$ 0).

Il valore non è un'asserzione di correttezza di \texttt{BrokerImpl}: è una stima statistica
relativa a questo campione di esecuzioni, al profilo operazionale uniforme assunto e
ad una suite che esercita 11 dei circa 250 metodi della classe, overload inclusi: 9 metodi
distinti, di cui \texttt{lock} e \texttt{detachAll} con due overload ciascuno.
